\documentclass[landscape,twocolumn,leqno,a4paper,10pt]{report} %leqno affiche les numéros d'équations à gauche

\usepackage[a4paper,top=0.4cm,bottom=0.95cm,left=0.4cm,right=0.4cm,marginparwidth=0.5cm,footskip=0.7cm]{geometry}



%----------------- Codage de la langue et de la police--------------------------------------------------------------------------
\usepackage[french]{babel} 		% Permet de spécifier au compilateur que l'on désire taper son document en français.
\usepackage[utf8x]{inputenc} 	% Permet de taper directement à l'écran, dans le fichier .tex, les caractères accentués.
\usepackage[T1]{fontenc} 		% Permet juste d'utiliser une nouvelle norme LATEX concernant le codage des caractères.
%\usepackage{lmodern}		% Pack de police --> Pas besoin puisque j'utilise Fourier

%------------------------ Packages utiles-------------------------------------------------------------------------------------------

\usepackage{graphicx}	% Pour insérer des images
\usepackage{caption}		% Légende

\newcommand\hmmax{0}		% Alors c'est une astuce assez mystérieuse pour dépasser la limite de polices latex.
							% En effet, le package bm est lourd et avec tout ce que l'on met, on obtient un message d'erreur : Too many math alphabets used in version normal.			
\newcommand\bmmax{0}		% Cette astuce règle le problème, il faut tout de même mettre ces compteurs avant le package bm

\usepackage{amsmath,amssymb}	% Permettent d'obtenir des caractères mathématiques (de toute beauté !).
\usepackage{amsthm}		% Pour pouvoir faire les newthm
\usepackage{mathtools}	% Corrige des bugs dans amsmath et le complète
\usepackage{numprint}	% Pour les espacements dans les nombres : \numprint{10000} = 10 000
\usepackage{fourier}		% REMARQUE : [upright] impose aux majuscules d'être droite dans le mode mathématique 
\usepackage{pifont} 		% Importe une large gamme de symboles, on utilise \ding{numéro}
\usepackage{xargs} 		% Permet de définir des macros à plusieurs paramètres optionnels.
\usepackage{fancyhdr} 	% Package pour les entêtes et bas de page
\pagestyle{fancy} 			% Accompagne \usepackage{fancyhdr}
\usepackage{bm} 		% Permet de mettre en gras du contenu mathématiques via \bm{..}
\usepackage{mathrsfs} 	% Pour les lettres cursives
\usepackage{enumitem} 	% Pour pouvoir personnaliser les listes

\frenchbsetup{StandardLists=true} % Afin d'éviter tout conflit entre les extensions enumitem et french babel. 

\usepackage{sectsty}	% Pour colorer les titres
\usepackage{titlesec} 	% Pour personnaliser les titres
\usepackage{setspace} 	% Pour changer la taille des interlignes
       	               
\usepackage{array,multirow,makecell}
\usepackage{hlist}		% Package Christian Tellechea sur les listes numérotées gauche vers droite
\setdefaulthlist{left margin = 2mm, label sep= 0.2mm}	% réglages package hlist
				

%------------ Les packages de dessin sous pstricks : il faut changer la compilation de latex pour compiler ces dessins------------------
% Options\Configurer Texmaker\Compil rapide\ Latex+dvips+ps2pdf+Voir PDF

\usepackage[pdf]{pstricks} 	% Pour faire des dessins
\usepackage{pstricks-add,pst-plot,pst-tree,pst-node,pst-math,pst-eucl,pst-func,pst-3dplot}% ,pst-text,pst-eps,pst-fill,pst-grad
\usepackage{multido}		% Pour faire des multido
\usepackage{lastpage} 	% Pour le compteur de la dernière page, on appelle alors celle-ci avec lastage



\usepackage[colorlinks=true, 
			%allcolors=blue
			]{hyperref}		% ATTENTION : doit être importé après la majorité des packages sauf tabularx et breakrules (par exemple)
\hypersetup{			% Paramétrage des hyperliens
	colorlinks = true,   		% Colorise les liens
	breaklinks = true,		% Permet les retours à la ligne pour les liens trop longs
	urlcolor = blue,		% Couleur des hyperliens
	linkcolor = black,		% Couleur des liens internes aux documents (index, figures, tableaux, équations, ...)
	citecolor = blue		% Couleur des liens vers les références bibliographiques
	}




%%%%%%%%%%%%%%%%%%%%%%%%%%%%%%%%%%%%%%%%%%%%%%%%%%
%------------------------- Définition des commandes -----------------------------------------------------------------------------
%%%%%%%%%%%%%%%%%%%%%%%%%%%%%%%%%%%%%%%%%%%%%%%%%%


%--------------------- Raccourcis fontes--------------------------------------------------------------
% REMARQUE : le package fourier met la même police pour mathcal et mathscr, on remet donc mathcal normalement et on utilise la définition de fourier (un peu plus grasse) pour mathscr
\DeclareMathAlphabet{\mathcal}{OMS}{cmsy}{m}{n}	
\DeclareMathAlphabet{\mathscr}{FMS}{futm}{m}{n}

%---------------------------- Raccourcis des espaces--------------------------------------------------------------------------------------------
\newcommand\n{\noindent}
\newcommand\q{\quad}
\newcommand\qq{\quad\quad}

%--------------------- Raccourcis en mode mathématiques-------------------------------------------------------------------------------------
\newcommand{\et}{\text{ et }}
\newcommand{\ou}{\text{ ou }}
\newcommand{\SSI}{\Longleftrightarrow}
\let\si\relax
\newcommand{\si}{\text{ si }}
\newcommand{\sinon}{\text{ sinon }}
\renewcommand{\epsilon}{\varepsilon}		% epsilon beaucoup plus joli
\renewcommand{\geq}{\geqslant} 
\renewcommand{\leq}{\leqslant}
\newcommand{\e}{\mathrm{e}}
\renewcommand{\textbf}[1]{\begingroup\bfseries\mathversion{bold}#1\endgroup}		% Très important, permet de passer en gras même avec le mode mathématiques

\newcommand{\Mod}[1]{\allowbreak\if@display\mkern5mu\else\mkern5mu\fi\left[#1\right]}	% Pour modulo entre crochets [.] avec un espacement convenable 


%---------------------------------- Raccourcis calcul diff----------------------------------------------------------------------------------------
\renewcommand{\d}{\;\mathrm{d}}	% Pour le d droit de l'intégrale car ce n'est pas une variable...


%-------------------------------- Raccourcis des ensembles----------------------------------------------------------------------------
\newcommand\R{\mathbb{R}}
\newcommand\C{\mathbb{C}}
\newcommand\K{\mathbb{K}}
\newcommand\Q{\mathbb{Q}}
\newcommand\D{\mathbb{D}}
\newcommand\Z{\mathbb{Z}}
\newcommand\N{\mathbb{N}}
\newcommand\U{\mathbb{U}}


%------------------------------- Dimensions intelignes, items...-----------------------------------------------------------------------------------------------
\onehalfspacing	% Pour que l'interligne soit de 1.5. Pour que l'interligne soit de la valeur que l'on souhaite, de base : \setstretch{1}

\setlist[1]{font=\bfseries,itemsep=0.2mm,leftmargin=6mm}
\setlist[2]{font=\bfseries,itemsep=0.2mm,leftmargin=4mm}
\setenumerate[1]{font=\bfseries,itemsep=0.2mm,leftmargin=6mm}
\setenumerate[2]{font=\bfseries, label=\alph*., leftmargin=4mm}
%%\setenumerate[profondeur]{font=\bfseries, label=\ding{221} \color{blue}} permet de changer en a) b) les sous-items d'une liste enumerate
%% itemsep = correspond à l'interligne entre chaque item
% de base : labelsep=2mm et permet de changer la distance entre l’étiquette et le texte des éléments
% leftmargin =4mm spécifie la marge à gauche de la liste
\AtBeginDocument{\renewcommand{\labelitemi}{\textbullet}}	% Pour que les puces par défaut soit des points.
%\AtBeginDocument{\renewcommand{\labelitemi}{\ding{86}}}	% Pour que les puces par défaut soit des étoiles.
\AtBeginDocument{\renewcommand{\labelitemii}{$\ast$}}		% Pour que les puces par défaut pour une sous-liste soit des astérisques.
\newcommand{\itemmm}{\item[$\ast$]}				% \itemmm devient l'item \ast
\newcommand{\itemmmm}{\item[\footnotesize$\circledcirc$\normalsize]}
\newcommand{\itemmmmm}{\item[\footnotesize$\astrosun$\normalsize]}
\newcommand{\itemmmmmm}{\item[\footnotesize$\multimap$\normalsize]}
\newcommand{\itemmmmmmmmm}{\item[$\divideontimes$]}
\newcommand{\etoile}{\item[$\star$]}
\newcommand{\etoilee}{\item[$\bigstar$]}
\newcommand{\losange}{\item[$\diamondsuit$]}
\newcommand{\carre}{\item[\LARGE$\centerdot$]}
\newcommand{\carreombre}{\item[\ding{112}]}

\setlength{\columnsep}{30pt} 		% Permet d'écarter les deux colonnes du twocolumn de 30pt, pratique pour écarter le texte de la ligne centrale
\setlength{\columnseprule}{0.5pt}	% Affiche une ligne verticale entre chaque item et entre chaque page du twocolumn
						% \columnseprule=0cm 	% Permet d'enlever le trait vertical entre chaque colonne
%\renewcommand\tabenumsep{0.5cm\hskiplem}	% Ecarte de 0.5cm les items de tabenum, sinon ils sont trop collés...


\newlength{\tagmarginsep} 		% On crée une longueur
\setlength{\tagmarginsep}{0cm}		% Permet de régler la position d'un tag pour une équation mettre avant l'équation : \setlength{\tagmarginsep}{x cm} pour que le tag se rapproche de l'équation
\everydisplay{\displayindent=\tagmarginsep \displaywidth=\dimexpr\linewidth-2\tagmarginsep}	% Va avec ce qu'il y a juste avant


%-------------------- Définition des couleurs-----------------------------------------------------------------------------------------

\definecolor{violetfonce}{RGB}{126,62,142}
\definecolor{magentafonce}{RGB}{128, 0, 128}
\newrgbcolor{magentafonce}{.5 0 .5}			% Une autre façon de définir une couleur, permet surtout d'utiliser \magenta
\definecolor{violet1}{cmyk}{0.5,1,0,0}
\definecolor{violet2}{cmyk}{0.24998, 1, 0, 0}
\definecolor{violet3}{cmyk}{0, 1, 0.25, 0}



\newcommand\tred{\textcolor{red}}
\newcommand\tredf{\textcolor{rougefonce}}
\newcommand\tdred{\textcolor{deepred}}
\newcommand\tblue{\textcolor{blue}}
\newcommand\tbluef{\textcolor{bleufonce}}
\newcommand\tgreen{\textcolor{vertfonce}}
\newcommand\torange{\textcolor{orange}}
\newcommand\tmagenta{\textcolor{magentafonce}}


%------------------ Changement paramètres des parties, sous parties---------------------------------------------------------


\sectionfont{\color{violet1}}			% Utilisation du package sectsty
\subsectionfont{\color{violet2}}
\subsubsectionfont{\color{violet3}}


%--------------------------------- Raccourcis probabilités--------------------------------------------------------------------------------
\newcommand{\Prob}{\mathbb{P}}
\newcommand{\E}{\mathbb{E}}
\newcommand{\normale}[2]{\mathcal{N}\left(#1;#2 \right)}	% Pour la loi Normale : 	$\normale{\mu}{\sigma^2}$
\newcommand{\Unif}[1]{\mathcal{U}\left(#1\right)}			% Pour la loi uniforme :	$\Unif{[a;b]}$
\newcommand{\Ber}[1]{\mathcal{B}\left(#1\right)}			% Pour la loi de Bernoulli : $\Ber{p}$
\newcommand{\Bino}[2]{\mathcal{B}\left(#1;#2\right)}		% Pour la loi Binomiale :	$\Bino{n}{p}$
\newcommand*{\binomial}[2]{% 
  \ensuremath{ 
    \begin{pmatrix} 
      #2\\ 
      #1
    \end{pmatrix}}}				% Pour le coefficient binomial : $\binom{k}{n}$
\renewcommand{\binom}{\binomial} 	% La commande \binom existait déjà mais n'était pas pratique



%----------------------------------- Opérateurs---------------------------------------------------------------------------------------------------

\newcommand{\barre}[1]{\overline{#1}}  									% Barre au dessus du texte, pratique pour le conjugué des complexes
\newcommandx{\Sum}[3][3=\white]{{\displaystyle\sum\limits_{\substack{#1 #3}}^{#2}}}		% $\Sum{k=0}{n}[<option>]$ où [<option>] est placée à la suite de k=0 si l'on tape de façon classique, et en-dessous si l'on met [\\ k\neq k_0]
															% On peut juste taper : $\Sum{k=0}{n}$ et cela fonctionne
															% \displaystyle impose de gérer les tailles comme en mode mathématique : ici on aura donc toujours de grande somme que ce soit en math comme dans le texte
\newcommandx{\Prod}[3][3=\white]{{\displaystyle\prod\limits_{\substack{#1 #3}}^{#2}}}		% $\Prod{k=0}{n}[<option>]$	idem pour l'option
\newcommand{\Int}[4]{{\displaystyle\int_{#1}^{#2} #3  \d #4}}   					% $\Int{0}{\pi}{\cos(t)}{t}$
\newcommandx{\Lim}[3][3=\white]{\lim\limits_{\substack{#1 \to #2 #3}}}		% $\Lim{x}{0}[<option>]$	idem pour l'option


%------------------------------------ Délimiteurs--------------------------------------------------------------------------------------------------
%\newcommand{\abs}[1]{\left\lvert #1 \right\rvert}			% Valeur absolue | |
%\newcommand{\norm}[1]{\left\lVert #1 \right\rVert}			% Doubles barres verticales : norme d'un vecteur || ||

%--------------------------------------- Vecteurs----------------------------------------------------------------------------------------------------
 
\newcommand{\Oij}{\left(O;\vec{\imath},\vec{\jmath}\right)}
\newcommand{\Oijk}{\left(O;\vec{\imath},\vec{\jmath},\vec{k}\right)}
\newcommand{\Ouv}{\left(O;\vec{u},\vec{v}\right)}
\newcommand{\Ouvw}{\left(O;\vec{u},\vec{v},\vec{w}\right)}

%------------------------------------ Systèmes d'équations-------------------------------------------------------------------------------------------
% Remarque \syst et \syste ne font pas d'alignement mais ils sont bien, si l'on veut de l'alignement, on utilisera \systeme du package "systeme"
\newcommand*{\syst}[2]{%			% Systèmes 2 éq. à ? inconnues : 	\syst{ax+by+c=0}{a'x+b'y+c'=0}
  \left\{\begin{array}{ll} 
        #1\\ 
        #2 
      \end{array}  
    \right. }
 
\newcommand*{\syste}[3]{% 			% Systèmes 3 éq. à ? inconnues : 	\syste{ax+by+c=0}{a'x+b'y+c'=0}{a''x+b''y+c''=0}
  \left\{\begin{array}{ll} 
        #1\\ 
        #2\\
        #3 
      \end{array}  
    \right.}



%------------------------------------------ Autre----------------------------------------------------------------------------------------------------------
\newcommand{\boite}[2]{\noindent\parbox[]{#1 cm}{#2} }  									% Boite autre façon de faire minipage en gros


%---------------------------- En-têtes et bas de pages----------------------------------------------------------------------------
% A chaque fois : *head* pour le haut et *foot* pour le bas
% L = left, C = center, R = right 
% L = toutes les pages, LE = toutes les pages paires, LO = toutes les pages impaires, idem CE, CO, ...


\newcommand{\fcours}[1]{								% f pour fancy				   
	\pagestyle{fancy}									% Efface le style qui a pu être mis avant
	\renewcommand{\headrulewidth}{0pt} 
	\renewcommand{\footrulewidth}{0.5pt} 
	\renewcommand{\footruleskip}{-1pt}						% Ecart vertical entre la ligne du bas et le texte en bas de page (texte foot)
	%\fancyhead[L]{}									
	%\fancyhead[C]{}
	%\fancyhead[R]{}									
	\fancyfoot[L]{\footnotesize \textit{#1}} 			% Classe : <Seconde> 
	\fancyfoot[C]{\footnotesize \thepage/\pageref{LastPage}}
	\fancyfoot[R]{\textit{\footnotesize C. Marguerite}}	% N'apparaîtra pas si on met les marges avec moins de bottom=1.3cm dans le package geometry
 }												% Fin de la commande
					%% Mettre au début du document : \fcours{Seconde} par exemple
					%% \usepackage[a4paper,top=0.4cm,bottom=1cm,left=0.4cm,right=0.4cm,marginparwidth=0.5cm,footskip=0.6cm]{geometry} 




%---------------------------- Réglages chapitres + table des matières------------------------------------------------------
\titleformat{\part}[display]
{\centering\normalfont\Huge\bfseries}
{\partname}
{20pt}
{\rule{6cm}{1pt}\raisebox{-2.1pt}[10pt][10pt]{\quad\decofourleft\decoone\decofourright\quad}\rule{6cm}{1pt} ~\\\vspace{20pt}}

\makeatletter
\renewcommand\chapter{\global\@topnum\z@
                    \@afterindentfalse
                    \secdef\@chapter\@schapter}
\def\@topnewpage[#1]{#1}
\makeatother										% Pour avoir le chapitre sur une demie page

\titleformat{\chapter}{\centering\normalfont\Large\bfseries\scshape\color{magentafonce}}{}{1em}{}
\titlespacing*{\chapter}{0pt}{-18pt plus 0pt minus 2pt}{0pt plus 2pt minus 2pt}		% Réglages espacements (à gauche, au-dessus, en-dessous) de la commande \chapter

\makeatletter
\newcommand{\twocolumntoc}{		
  \section*{\huge\centering\contentsname
    \@mkboth{%
      \MakeUppercase\contentsname}{\MakeUppercase\contentsname}}%
  \@starttoc{toc}%
}
\makeatother										% Permet de mettre la table des matières en twocolumn 
												% Mettre après \begin{doc..} \twocolumntoc \clearpage 		% Remplace \tableofcontents par \twocolumntoc en twocolumn

\makeatletter
\renewcommand{\thesection}{\@arabic\c@section}
\makeatother										% Enlève la numérotation des chapitres devant les sections

\setcounter{tocdepth}{3} 								% Permet d'afficher dans la table des matières suivant une certaine profondeur : -1=part, 0=chapter, 1=section, ...
\setcounter{secnumdepth}{3}							% Pour avoir la numérotation des subsubsections avec les classes book et report

\newcommand{\nocontentsline}[3]{}
\let\oldaddcontentsline\addcontentsline
\newcommand{\tocless}[2]{\let\addcontentsline\nocontentsline #1{#2}\let\addcontentsline\oldaddcontentsline}	% Pour retirer une section de la table des matières : \tocless\section{...}

%Tapuscrit : Christian Marguerite
\setlength{\parindent}{0pt}


\setstretch{1}
\begin{document}
\fcours{Concours général}

\chapter{Concours général des lycées - Session 2023}

\section{Problème 1 : Soyons rationnels !}


Pour tout entier $n \geq 1$, on note $v(n)$ le plus grand entier $k$ tel que $\frac{n}{2^{k}}$ soit un entier.\\
On définit la suite $\left(u_{n}\right)_{n \geq 1}$ par récurrence, en posant $u_{1}=1$ puis, pour tout entier $n \geq 2$,

$$u_{n}= \begin{cases}
0 & \text { si } u_{n-1}=0 \\ 
1+2 v(n)-\frac{1}{u_{n-1}} & \text { si } u_{n-1} \neq 0\end{cases}$$

\begin{enumerate}
	\item Donner la valeur des entiers $v(1)$, $v(2)$, $v(3)$ et $v(4)$.
	\item Démontrer, pour tout entier $n \geq 1$, que $v(n)=0$ si $n$ est impair et que $v(n)=v\left(\frac{n}{2}\right)+1$ si $n$ est pair.
	\item Calculer les huit premiers termes de la suite $\left(u_{n}\right)_{n \geq 1}$ et vérifier que $u_{8}=4$.
	\item Démontrer, pour tout entier $n \geq 1$, que $u_{n}$ est un nombre rationnel strictement positif, que $u_{2 n}=u_{n}+1$ et que $u_{2 n+1}=\frac{u_{n}}{u_{n}+1}$.
	\item Démontrer que tout nombre rationnel strictement positif est égal à un terme $u_{n}$.
	\item Démontrer que tout nombre rationnel strictement positif est égal à un unique terme $u_{n}$.
\end{enumerate}

\section{Problème 2 : Limite sympathique !}

\subsection{Quelques exemples}


\begin{enumerate}
	\item On considère dans cette question, pour tout entier $n \geq 1$, l'équation :
$$x^{2}+\frac{1}{n} x-1=0$$
d'inconnue $x$.
	\begin{enumerate}
		\item Soit $n$ un entier naturel non nul. Démontrer que cette équation admet une unique solution réelle positive; on la note $x_{n}$. Exprimer $x_{n}$ en fonction de $n$.
		\item Démontrer que la suite $\left(x_{n}\right)_{n \geq 1}$ converge, on note $x_{\infty}$ sa limite.
		\item Démontrer que $x_{\infty}$ est solution de l'équation $x^{2}-1=0$.
	\end{enumerate}
	\item On considère dans cette question, pour tout entier $n \geq 1$, l'équation :
	$$\frac{1}{n} y^{2}-y-1=0$$
	d'inconnue $y$.
	\begin{enumerate}
		\item Soit $n$ un entier naturel non nul. Démontrer que cette équation admet une unique solution réelle positive, on la note $y_{n}$.
		\item Démontrer que la suite $\left(y_{n}\right)_{n \geq 1}$ diverge.
	\end{enumerate}
	\item On considère dans cette question, pour tout entier $n \geq 1$, l'équation :
	$$z^{3}+\frac{1}{n} z^{2}-1=0$$
	d'inconnue $z$.
	\begin{enumerate}
		\item Soit $n$ un entier naturel non nul.
		\begin{enumerate}
			\item Étudier les variations de la fonction $z\mapsto z^3+\frac{1}{n}z^2-1$ sur l'intervalle $[0,+\infty[$.
			\item En déduire que cette équation admet une unique solution réelle positive, on la note $z_{n}$. Démontrer que $z_{n}$ appartient à l'intervalle $]0,1[$.
		\end{enumerate}
		\item Démontrer que la suite $\left(z_{n}\right)_{n \geq 1}$ est convergente.\\
		\textit{On pourra s'intéresser au signe du réel \quad $z_{n+1}^{3}+\frac{1}{n} z_{n+1}^{2}-1$.}
		\item On note $z_{\infty}$ la limite de la suite $\left(z_{n}\right)_{n \geq 1}$.\\
		Démontrer que $z_{\infty}$ est solution de l'équation $z^{3}-1=0$.
	\end{enumerate}
	\item On considère dans cette question, pour tout entier $n \geq 1$, l'équation :
	$$\frac{1}{n} t^{3}-t^{2}-1=0$$
	d'inconnue $t$.
	\begin{enumerate}
		\item Soit $n$ un entier naturel non nul. Démontrer que cette équation admet une unique solution réelle, on la note $t_{n}$.
		\item La suite $\left(t_{n}\right)_{n \geq 1}$ est-elle convergente ? Si oui, quelle est sa limite ?
	\end{enumerate}
\end{enumerate}


\subsection{Polynômes sympathiques}


Dans les deux prochaines parties, on considère un entier $d \geq 1$. La fonction $P$ est un \textit{polynôme de degré au plus $d$} s'il existe des réels $a_{0}, a_{1}, \ldots, a_{d}$ tels que :
$$P(x)=a_{d} x^{d}+a_{d-1} x^{d-1}+\cdots+a_{2} x^{2}+a_{1} x+a_{0}$$
pour tout réel $x$.\\
Soit $P: x \mapsto a_{d} x^{d}+a_{d-1} x^{d-1}+\cdots+a_{2} x^{2}+a_{1} x+a_{0}$ un polynôme de degré au plus $d$. On dit que :

\begin{itemize}
	\item[$\triangleright$] $P$ est \textit{initialement sympathique} si $a_{0}=-1$ et si $a_{k} \geq 0$ pour tout entier $k$ tel que $1 \leqslant k \leqslant d$;
	\item[$\triangleright$] $P$ est \textit{faussement sympathique} si $a_{0}=-1$ et si $a_{k} \leqslant 0$ pour tout entier $k$ tel que $1 \leqslant k \leqslant d$;
	\item[$\triangleright$] $P$ est \textit{vraiment sympathique} si $a_{0}=-1$ et s'il existe un entier $k$ tel que $1 \leqslant k \leqslant d-1$ et pour lequel $a_{1} \leqslant 0, a_{2} \leqslant 0, \ldots, a_{k} \leqslant 0$ et $a_{k+1}>0, a_{k+2} \geq 0, \ldots, a_{d} \geq 0$.
\end{itemize}
Enfin, on dit que $P$ est \textit{sympathique} s'il est initialement, faussement ou vraiment sympathique.

\begin{enumerate}\setcounter{enumi}{4}
	\item Quels sont les polynômes qui sont à la fois faussement sympathiques et initialement sympathiques ?
	\item Démontrer que tout polynôme faussement sympathique est
	\begin{enumerate}
		\item strictement négatif sur l'intervalle $[0,+\infty[$;
		\item décroissant sur l'intervalle $[0,+\infty[$.
	\end{enumerate}
	\item Soit $P$ un polynôme vraiment sympathique et initialement sympathique.
	\begin{enumerate}
		\item Démontrer que $P$ est strictement croissant sur l'intervalle $[0,+\infty[$;
		\item Démontrer que l'équation $P(x)=0$ admet une unique solution strictement positive.
	\end{enumerate}
	\item Soit $P$ un polynôme vraiment sympathique mais pas initialement sympathique.
	\begin{enumerate}
		\item Démontrer qu'il existe un réel $b>0$, un entier $\ell \geq 0$ et un polynôme $Q$ vraiment sympathique tels que :
		$$P'(x)=b x^{\ell} Q(x)$$
pour tout réel $x$.
		\item Démontrer qu'il existe un réel $r>0$ tel que le polynôme $P$ vérifie les quatre propriétés suivantes :
		\begin{itemize}
			\item[$\triangleright$] $P$ est décroissant sur l'intervalle $[0, r]$;
			\item[$\triangleright$] $P$ est strictement croissant sur l'intervalle $[r,+\infty[$;
			\item[$\triangleright$] $P$ est strictement négatif sur l'intervalle $[0, r]$;
			\item[$\triangleright$] l'équation $P(x)=0$ admet une unique solution dans l'intervalle $[r,+\infty[$.
		\end{itemize}
	\end{enumerate}
	\item Quels sont les polynômes sympathiques $P$ pour lesquels l'équation $P(x)=0$ admet au moins une solution strictement positive ? Donner, dans ce cas, le tableau de signes de $P$ sur l'intervalle $[0,+\infty[$.
\end{enumerate}


\subsection{De la suite dans les idées}

On considère désormais des polynômes vraiment sympathiques $P_{1}, P_{2}, \ldots$ Puisque ces polynômes sont de degré au plus $d$, on peut écrire chaque polynôme $P_{n}$ sous la forme :
$$P_{n}: x \mapsto a_{d, n} x^{d}+a_{d-1, n} x^{d-1}+\cdots+a_{2, n} x^{2}+a_{1, n} x+a_{0, n}$$
On suppose en outre, pour tout entier $k$ tel que $0 \leqslant k \leqslant d$, que la suite $\left(a_{k, n}\right)_{n \geq 1}$ est convergente, on note $a_{k, \infty}$ sa limite.\\
On considère alors le polynôme $P_{\infty}$ défini par :
$$P_{\infty}: x \mapsto a_{d, \infty} x^{d}+a_{d-1, \infty} x^{d-1}+\cdots+a_{2, \infty} x^{2}+a_{1, \infty} x+a_{0, \infty}$$
Enfin, pour tout entier $n \geq 1$, on note $x_{n}$ l'unique solution strictement positive de l'équation $P_{n}(x)=0$. Ci-dessous, on étudie la convergence éventuelle de la suite $\left(x_{n}\right)_{n \geq 1}$.

\begin{enumerate}\setcounter{enumi}{9}
	\item Soit $t$ un réel fixé. Démontrer que la suite $\left(P_{n}(t)\right)_{n \geq 1}$ converge vers $P_{\infty}(t)$.
	\item Démontrer que le polynôme $P_{\infty}$ est sympathique.
	\item On suppose dans cette question que le polynôme $P_{\infty}$ est vraiment sympathique, et on note $x_{\infty}$ l'unique solution strictement positive de l'équation $P_{\infty}(x)=0$.
	\begin{enumerate}
		\item Soient $u$ et $v$ deux réels tels que $0<u<x_{\infty}<v$. Démontrer qu'il existe un entier $M_{u, v}$ tel que $P_{n}(u)<0<P_{n}(v)$ pour tout entier $n \geq M_{u, v}$.
		\item En déduire que la suite $\left(x_{n}\right)_{n \geq 1}$ converge vers $x_{\infty}$.
	\end{enumerate}
	\item On suppose dans cette question que le polynôme $P_{\infty}$ est faussement sympathique. Démontrer que $\left(x_{n}\right)_{n \geq 1}$ diverge vers $+\infty$.
	\item Retrouver les résultats de la première partie.
\end{enumerate}



\section{Problème 3 : Polynômes et polygones réguliers}

Le plan est rapporté à un repère orthonormé $\mathscr{R}=(O ; \vec{\imath}, \vec{\jmath})$.\\
Soit $k$ un entier tel que $k \geq 3$. Les points $M_{1}, M_{2}, \ldots, M_{k}$ sont les sommets d'un polygone régulier de centre $O$ si ces points
\begin{itemize}
	\item[$\triangleright$] sont deux à deux distincts,
	\item[$\triangleright$] apparaissent dans le sens trigonométrique (c'est-à-dire le sens contraire des aiguilles d'une montre) sur un même cercle de centre $O$, et
	\item[$\triangleright$] vérifient l'égalité $M_{1} M_{2}=M_{2} M_{3}=\cdots=M_{k-1} M_{k}=M_{k} M_{1}$. En particulier, pour $k=3$, il s'agit d'un triangle équilatéral; pour $k=4$, il s'agit d'un carré.
\end{itemize}
Pour tout entier $d \geq 0$, une fonction $P$ est un polynôme de degré $d$ s'il existe des réels $a_{0}, a_{1}, \ldots, a_{d}$ tels que $a_{d} \neq 0$ et
$$P(x)=a_{d} x^{d}+a_{d-1} x^{d-1}+\cdots+a_{2} x^{2}+a_{1} x+a_{0}$$
pour tout réel $x$; on pourra admettre que, pour un tel polynôme, l'équation $P(x)=0$ admet au plus $d$ solutions réelles.\\
Quant à elle, la fonction $P: x \mapsto 0$ est appelée le polynôme nul.\\
Enfin, étant donné un polynôme $P$ (nul ou non), on note $\mathscr{C}_{P}$ la courbe représentative de $P$ dans le repère $\mathscr{R}$.

\subsection{Triangles équilatéraux}

\begin{enumerate}
	\item Soit $P$ un polynôme de degré 1. Existe-t-il un triangle équilatéral dont les sommets appartiennent à $\mathscr{C}_{P}$ ?
	\item On considère les points :
	$$A\left(1, \frac{\sqrt{3}}{3}\right)\quad,\quad B\left(-1, \frac{\sqrt{3}}{3}\right) \quad\text { et }\quad C\left(0,-\frac{2 \sqrt{3}}{3}\right)$$
	\begin{enumerate}
		\item Démontrer que $A$, $B$ et $C$ sont les sommets d'un triangle équilatéral de centre $O$.
		\item Démontrer que les points $A$, $B$ et $C$ appartiennent à la courbe représentative du polynôme :
		$$Q: x \mapsto \frac{\sqrt{3}}{3}\left(3 x^{2}-2\right)$$
		\item Démontrer que les points $A$, $B$ et $C$ appartiennent à la courbe représentative du polynôme :
		$$R: x \mapsto \frac{\sqrt{3}}{3}\left(3 x^{2}-2\right)+x\left(x^{2}-1\right)$$
		\item Démontrer que, pour tout entier $d \geq 2$, il existe un polynôme de degré $d$ dont la courbe représentative contient les points $A$, $B$ et $C$.
	\end{enumerate}
\end{enumerate}

\subsection{Carrés de centre $O$}


Dans les questions \textbf{3.} et \textbf{4.}, on considère un polynôme $P$ et un carré $ABCD$ de centre $O$ dont les quatre sommets appartiennent à $\mathscr{C}_{P}$.

\begin{enumerate}\setcounter{enumi}{2}
	\item \begin{enumerate}
		\item Exprimer les coordonnées des points $B$, $C$ et $D$ en fonction de celles de $A$. Démontrer que les abscisses de $A$, $B$, $C$ et $D$ sont distinctes et non nulles.
		\item Démontrer que $P$ est non nul et que son degré vaut au moins 3.
	\end{enumerate}
	\item On suppose dans cette question qu'il existe des réels $a$, $b$ et $c$ tels que :
	$$P: x \mapsto x^{3}+a x^{2}+b x+c$$
	\begin{enumerate}
		\item Démontrer que $a=0$ et $c=0$.
		\item Démontrer que les abscisses respectives de $A$, $B$, $C$ et $D$ sont solutions de l'équation :
		$$P(P(x))+x=0$$
		\item Démontrer que le polynôme :
		$$Q: x \mapsto x^{4}+3 b x^{3}+3 b^{2} x^{2}+b\left(b^{2}+1\right) x+b^{2}+1$$
		admet au moins deux racines positives distinctes.
		\item Démontrer que $b<0$.
		\item On suppose qu'il existe deux réels $\alpha$ et $\beta$ tels que $0<\alpha<\beta$ et
		$$Q(x)=(x-\alpha)^{2}(x-\beta)^{2}$$
		pour tout réel $x$. Démontrer qu'alors $b=-\sqrt{8}$, puis déterminer les valeurs de $\alpha$ et $\beta$.
	\end{enumerate}
	\item \begin{enumerate}
		\item Démontrer qu'il existe un polynôme $P$ de degré 3 et un carré $ABCD$ de centre $O$ dont les sommets appartiennent à $\mathscr{C}_{P}$.
		\item Pour quels entiers $d$ existe-t-il un polynôme de degré $d$ dont la courbe représentative contient les points $A$, $B$, $C$ et $D$ obtenus en question \textbf{5.a.} ?
	\end{enumerate}
\end{enumerate}


\subsection{Où l'on prouve que $d \geq k-1$}

Soit $M_{1} M_{2} \cdots M_{k}$ un polygone régulier de centre $O$. On suppose dans cette question qu'il existe un polynôme $P$, de degré $d$, dont la courbe contient les points $M_{1}, M_{2}, \ldots, M_{k}$. On souhaite alors démontrer que $d \geq k-1$.\\
Pour tout $i$, on note $\left(x_{i}, y_{i}\right)$ les coordonnées de $M_{i}$ dans le repère $\mathscr{R}$.

\begin{enumerate}\setcounter{enumi}{5}
	\item \begin{enumerate}
		\item Pourquoi peut-on supposer que $x_{1}$ est inférieur ou égal à $x_{2}, x_{3}, \ldots, x_{k}$ et que $y_{1} \leqslant 0$ ?
		\item Démontrer que les abscisses $x_{i}$ sont deux à deux distinctes et que les ordonnées $y_{i}$ sont non nulles.
		\item Démontrer qu'il existe un nombre réel $R>0$ et un nombre réel $\theta$ appartenant à l'intervalle $]0, \pi / k[$  tels que $x_{1}=-R \cos (\theta)$ et $y_{1}=-R \sin (\theta)$.
		\item Démontrer que $x_{1}<x_{k}<x_{2}<x_{k-1}<x_{3}<x_{k-2}<\cdots$
		\item Démontrer que $P$ admet une racine sur chacun des $k-1$ intervalles :
		$$]x_{1},x_{k}[ \quad ,\quad ]x_{k},x_{2}[ \quad,\quad ]x_{2},x_{k-1}[ \quad,\quad ]x_{k-1},x_{3}[ \quad,\quad ]x_{3},x_{k-2}[ \quad, \ldots$$
		\item En conclure que $d \geq k-1$.
	\end{enumerate}
\end{enumerate}


\subsection{Où l'on prouve que tout entier $d \geq k-1$ convient}


On suppose dans cette partie que les abscisses $x_{i}$ sont deux à deux distinctes et on veut démontrer que, pour tout entier $d \geq k-1$, il existe un polynôme de degré $d$ dont la courbe contient les points $M_{1}, M_{2}, \ldots, M_{k}$.

\begin{enumerate}\setcounter{enumi}{6}
	\item Soient $a$ et $b$ deux réels. Dans le repère $\mathscr{R}$, on considère les points :
	$$A(\cos (a), \sin (a)) \quad,\quad B(\cos (a+b), \sin (a+b)) \quad\text { et }\quad C(-\sin (a), \cos (a))$$
	\begin{enumerate}
		\item Démontrer que le repère $\mathscr{R}'=(O ; \overrightarrow{O A}, \overrightarrow{O C})$ est orthonormé.
		\item Quelles sont les coordonnées du point $B$ dans le repère $\mathscr{R}'$ ?
		\item En déduire que :
		$$\begin{aligned}
& \cos (a+b)=\cos (a) \cos (b)-\sin (a) \sin (b) ; \\
& \cos (a-b)=\cos (a) \cos (b)+\sin (a) \sin (b) .
\end{aligned}$$
	\end{enumerate}
	\item On considère la suite de polynômes définie par \quad $T_{0}: x \mapsto 1$,\quad $T_{1}: x \mapsto x$ \quad et 
	$$T_{n+2}: x \mapsto 2 x T_{n+1}(x)-T_{n}(x)$$
	pour tout entier $n \geq 0$.
	\begin{enumerate}
		\item Démontrer que $T_{n}(\cos (\theta))=\cos (n \theta)$ pour tout entier $n \geq 0$ et tout réel $\theta$.
		\item Soit $\theta$ un réel, et soient $\ell \geq 1$ et $j \geq 0$ deux entiers. Démontrer que :
		$$T_{\ell-1}\left(\cos \left(\theta+\frac{2 j \pi}{\ell}\right)\right)-\cos (\ell \theta) \cos \left(\theta+\frac{2 j \pi}{\ell}\right)=\sin (\ell \theta) \sin \left(\theta+\frac{2 j \pi}{\ell}\right)$$
		\item Démontrer que, pour tout entier $d \geq k-1$, il existe un polynôme de degré $d$ dont la courbe contient les points $M_{1}, M_{2}, \ldots, M_{k}$.
	\end{enumerate}
\end{enumerate}

\end{document}